% !TeX root = ../深度学习与大语言模型技术原理.tex
% 由主文档通过 \input 载入；请编译主文档。
\setcounter{chapter}{28} % 保留原书章号。
\chapter{工程工具链与可复现实验}
\label{ch:engineering}

\section{模型、数据与后训练框架的职责}

PyTorch 负责张量计算与自动微分；Hugging Face 生态中的 Transformers 提供模型结构、权重与生成接口，Tokenizers 负责分词实现，Datasets 组织数据读取与变换，PEFT 管理适配参数。TRL 组织若干后训练目标，verl 与 OpenRLHF 面向包含采样和分布式组件的后训练流程。理解这些职责，有助于把问题定位到数据、目标、模型或系统层，而不是在多个框架间反复试默认参数。

高层 Trainer 能封装训练循环，但有效词元归一化、梯度累积、模板、packing 和终止标记仍需核对。训练器升级可能改变默认设置，配置文件也未必覆盖所有隐式默认值。因此实验应保存解析后的完整配置、依赖锁定信息和代码版本，关键损失最好能在小样本上独立重算。

TensorBoard、Weights \& Biases 等工具记录损失、学习率、梯度、显存和任务指标。日志应区分训练步、优化器步、生成轮次与累计词元，避免不同横轴造成误读。只记录最终分数不足以分析训练何时开始退化，也不足以恢复被中断实验。

\section{Python、C++ 与系统基础}

Python 的类型标注和数据类帮助表达配置与接口，但运行时输入仍需验证。多进程适于并行数据处理或独立工作者，异步 I/O 适于等待网络工具；异步语法本身不会让 GPU 计算自动并行。队列需要背压（backpressure），防止生产速度长期超过消费速度导致内存增长。

C++ 基础涉及值、指针、引用、对象生命周期和线程同步。张量跨 Python 与原生扩展边界时，要明确谁拥有内存、何时释放以及设备异步操作是否结束。数据竞争、越界和悬空指针可能表现为偶发训练错误，因此底层加速前应先建立可靠接口与错误检查。

Linux 操作能力包括进程、日志、文件权限、磁盘与 GPU 监测，以及通过远程会话管理长任务。文件搜索可用 \texttt{rg}，日志和数据处理可结合 \texttt{awk}、\texttt{sed}，设备状态可用 \texttt{nvidia-smi}；工具名称只是入口，关键是识别任务是否受 CPU、I/O、显存、算力或通信限制。

\section{版本管理、容器与恢复}

Git 分支保存并行改动，合并（merge）保留分支关系，变基（rebase）重写提交基底，拣选（cherry-pick）移植特定提交。实验结果应关联确切提交与未提交改动状态；仅记录分支名不足以确定当时程序。冲突解决需要理解两侧意图，不能自动选择整份文件的一侧而假定正确。

Docker 镜像描述运行环境，容器是运行实例，卷（volume）或挂载保存持久数据，网络配置决定外部连接。应固定基础镜像与依赖版本，并明确 GPU 驱动和运行时兼容关系。把数据写在临时容器内部可能在删除后丢失；把宿主目录全部挂入又可能扩大访问范围。

可恢复训练需要模型、优化器、调度器、随机状态和数据进度相互一致。保存中断导致部分文件不完整时，应能识别检查点是否有效。验证恢复可比较同一小任务连续训练与中断恢复后的损失、步数和数据序列；分布式及非确定性内核下不一定逐位一致，应说明允许偏差。

\section{从失败样本到研究问题}

一个可解释实验先提出可证伪假设，例如“领域回答错误主要来自检索漏召回”，再选择能区分原因的干预，例如提供人工证据、替换召回器或扩大候选集。每次只改变必要因素，并固定数据和计算预算。若多个组件同时改变，最终分数提升也难以归因。

报告应交代任务、数据划分、基线、变量、预算、质量、成本、不确定性和失败类别。负结果同样有价值：某项技术没有带来收益，可能说明瓶颈不在该层，也可能说明实现或条件不满足。附录~\ref{app:experiments} 将清单中的六个项目组织为从公式、实现到评测证据的练习。

